Deux charges positives $q_1 = \qty{\v{q1}}{\micro\coulomb}$ et $q_2 = \qty{\v{q2}}{\micro\coulomb}$ sont fixées à \qty{\v{d}}{cm} l'une de l'autre. On veut placer une troisième charge $Q = \qty{\v{Q}}{nC}$ sur la droite qui les relie.

\begin{enumerate}[label=\alph*)]
    \item Sans calcul, expliquez pourquoi la position d'équilibre de $Q$ doit se trouver entre les deux charges.
    \item À quelle distance de $q_1$ faut-il placer $Q$ pour que la force résultante sur elle soit nulle ?
    \item Si on place plutôt $Q$ exactement au milieu des deux charges, quel est le module de la force résultante sur $Q$ ?
\end{enumerate}
